% =============================================================================
% appendix.tex — inputted by paper.tex after \appendix
% =============================================================================

\section{The PNADC and the advanced panel identification}\label{app:pnadc}

\paragraph{The survey.} The Continuous National Household Sample Survey (\emph{Pesquisa Nacional por Amostra de Domicílios Cont\'inua}, PNADC) is the household labor-force survey conducted by the Brazilian Institute of Geography and Statistics (IBGE). It is a stratified, multi-stage probability sample of the resident population, released quarterly, and is the source of Brazil's official labor-market statistics. The survey follows a rotating panel scheme: each selected household is interviewed in 5 consecutive quarters and then leaves the sample. As a result, a given household, and the individuals in it, can be observed up to 5 times over roughly 15 months, which provides the short longitudinal dimension exploited in this paper. Because the sample is not self-weighting, all estimates use the survey's calibrated sampling weights, and the design information is retained for variance estimation. The outcome used as the first stage, the location at which a worker performs her main job, has been collected since the first quarter of 2018, which sets the start of the estimation window.

\paragraph{Linking individuals across quarters.} The public microdata do not contain a permanent personal identifier, so respondents must be linked across quarters from the information that is available: the household's position in the sample rotation, the ordering of individuals within the household, and demographic characteristics such as sex and date of birth. When interviews are complete and consistent, this linkage is straightforward. In practice a non-trivial share of interviews contain small inconsistencies, such as a missing or misrecorded birth date or a change in the reported household roster, that break naive matches. I use an advanced identification procedure that reconstructs the panel in the presence of such inconsistencies by treating the linkage as a graph-matching problem, in which candidate records are connected when they are compatible on the stable identifiers and demographic characteristics and the most consistent set of links is selected. Individuals that the procedure cannot confidently match are assigned a singleton identifier and therefore contribute nothing to the within-person variation that identifies the fixed-effects estimates. The estimation sample is accordingly restricted to confidently matched individuals; the roughly 3.7\% of women who cannot be linked are excluded. Because the excluded observations are singletons, this restriction changes the number of observations but leaves the fixed-effects estimates numerically unchanged.

\paragraph{Panel quality.} Table~\ref{tab:panel_retention} documents that the reconstructed panel captures genuine repeated observations rather than spurious matches. It reports, for the estimation sample, the share of households and of matched individuals observed for at least a given number of quarters, together with the quarter-to-quarter continuation probability. The figures are computed on the analysis sample, prime-age women who head a household or are the head's spouse, and so reflect both survey attrition and the sample definition: a woman who ages out of the 18--49 window, or a household that no longer contains a qualifying woman, leaves the sample even though the survey continues to interview her. This is why retention is lower than an unrestricted population figure would be; it is a feature of the sample definition, not a failure of the matching. The key point for the identification is that a substantial share of women are observed repeatedly, so that individual fixed effects are meaningful.

\input{../analysis/output/tables/tabA0_panel_retention.tex}

\clearpage
\section{Additional tables and figures}\label{app:additional}

Table~\ref{tab:robustness} collects the robustness checks on the first stage discussed in Section~\ref{sec:robustness}, and Figure~\ref{fig:hetero} summarizes the subgroup estimates of Table~\ref{tab:heterogeneity} as a coefficient plot. Table~\ref{tab:did_outcomes_B} repeats the difference-in-differences estimates using the broad comparison group (women with no child aged 0--7), which delivers more precision at the cost of a less tightly matched counterfactual; the conclusions are unchanged. Figure~\ref{fig:event_maternity} shows the event study for maternity leave, which exhibits the pronounced pre-reform trend discussed in the text and motivates treating that outcome as non-causal. Table~\ref{tab:occ_transition} reports the occupation-transition matrix underlying the mechanism discussion. Table~\ref{tab:triple_diff_outcomes} extends the triple difference to all outcomes. Figure~\ref{fig:window} reports the control-window robustness.

\begin{table}[htbp]\centering
\caption{Robustness of the First-Stage Home-Office Effect}
\label{tab:robustness}\small
\begin{tabular}{lcc}
\toprule
 & Treated $\times$ Post & Obs. \\
\midrule
\multicolumn{3}{l}{\textit{Home office (pp)}} \\
Baseline & -0.08$^{}$ & 1,020,823 \\
 & (0.30) & \\
Alternative cutoff (2022Q1 post) & 0.20$^{}$ & 1,020,823 \\
 & (0.30) & \\
Treated excl.\ grandchildren & -0.07$^{}$ & 956,777 \\
 & (0.31) & \\
Treated excl.\ stepchildren & -0.13$^{}$ & 1,005,187 \\
 & (0.30) & \\
Excluding 2020--2021 (COVID) & -0.09$^{}$ & 823,117 \\
 & (0.37) & \\
Ages 20--35 & -0.36$^{}$ & 617,885 \\
 & (0.38) & \\
Ages 20--40 & -0.41$^{}$ & 830,421 \\
 & (0.33) & \\
State $\times$ quarter fixed effects & -0.08$^{}$ & 1,020,823 \\
 & (0.30) & \\
Two-way cluster (household, PSU) & -0.08$^{}$ & 1,020,823 \\
 & (0.29) & \\
Matched panel only & -0.08$^{}$ & 989,898 \\
 & (0.30) & \\
Telework-eligible only & -0.95$^{}$ & 139,496 \\
 & (1.01) & \\
\midrule
\multicolumn{3}{l}{\textit{Real income (R\$), outlier sensitivity}} \\
Income: raw & 32.7$^{}$ & 495,994 \\
 & (35.7) & \\
Income: winsor. 1\% & 5.3$^{}$ & 495,994 \\
 & (26.8) & \\
\bottomrule\end{tabular}
\par\vspace{3pt}\footnotesize\raggedright \textit{Notes:} Each estimate is a separate difference-in-differences regression on the preferred sample (treated vs.\ Control A, child 5--7), including the treated main effect, individual and year-quarter fixed effects, and survey weights; standard errors clustered at the household in parentheses (except the two-way row, clustered at the household and primary sampling unit). $^{*}$/$^{**}$/$^{***}$ denote significance at 10/5/1\%. The placebo on men is reported in Table \ref{tab:triple_diff}.
\end{table}


\begin{figure}[H]
\centering
\includegraphics[width=0.85\textwidth]{fig07_heterogeneity_coefplot.pdf}
\caption{First-stage home-office effect by subgroup. Points are difference-in-differences estimates and horizontal bars are 95\% confidence intervals from household-clustered standard errors. Subgroups are defined at baseline.}
\label{fig:hetero}
\end{figure}


\begin{table}[H]
   \caption{\label{tab:did_outcomes_B} Difference-in-Differences Estimates by Outcome, Control B}
   \centering
\footnotesize\setlength{\tabcolsep}{3pt}
   \begin{tabular}{lcccccc}
      \toprule
      Dependent Variables:     & Home office & Real income & Usual hours    & Employed       & In labor force & Maternity leave\\  
      Model:                   & (1)         & (2)         & (3)            & (4)            & (5)            & (6)\\  
      \midrule
      \emph{Variables}\\
      Treated (child $\leq$4)  & -0.003      & -20.3       & -0.645$^{***}$ & -0.065$^{***}$ & -0.071$^{***}$ & 0.114$^{***}$\\   
                               & (0.002)     & (36.4)      & (0.177)        & (0.004)        & (0.004)        & (0.004)\\   
\addlinespace[2pt]
      Treated $\times$ Post    & 0.004       & 14.3        & -0.093         & 0.008$^{*}$    & 0.008$^{**}$   & -0.011$^{***}$\\   
                               & (0.003)     & (39.5)      & (0.177)        & (0.004)        & (0.004)        & (0.002)\\   
\addlinespace[2pt]
      \midrule
      \emph{Fixed-effects}\\
      Individual               & Yes         & Yes         & Yes            & Yes            & Yes            & Yes\\  
      Year-quarter             & Yes         & Yes         & Yes            & Yes            & Yes            & Yes\\  
      \midrule
      \emph{Fit statistics}\\
      Observations             & 2,192,875   & 1,251,406   & 1,302,376      & 2,192,875      & 2,192,875      & 2,192,875\\  
      R$^2$                    & 0.664       & 0.889       & 0.765          & 0.797          & 0.778          & 0.475\\  
\bottomrule
   \end{tabular}
   
   \par \raggedright 
   \footnotesize\textit{Notes:} As in the Control~A outcome table, using the broad comparison group (women with no child aged 0--7). Each column is a separate difference-in-differences regression under the preferred specification (treated main effect with individual and year-quarter fixed effects, weighted by the survey weights). Home office, employed, in labor force, and maternity leave are 0/1 indicators, so a coefficient of $0.01$ corresponds to one percentage point; real income is in reais per month and usual hours in hours per week, both observed for workers only. Standard errors clustered at the household in parentheses. Significance levels: *** $p<0.01$, ** $p<0.05$, * $p<0.10$.
\end{table}




\begin{figure}[H]
\centering
\includegraphics[width=0.82\textwidth]{fig09_event_study_maternity.pdf}
\caption{Event study: maternity leave around the reform. Constructed exactly as Figure~\ref{fig:event_home}. The pronounced pre-reform trend in the treated-versus-control gap indicates that the parallel-trends assumption fails for this outcome, which is therefore not interpreted causally.}
\label{fig:event_maternity}
\end{figure}

\begin{table}[H]\centering
\caption{Occupation Transitions Around the Reform}
\label{tab:occ_transition}\small
\begin{tabular}{lcc}
\toprule
\multicolumn{3}{l}{\textit{Treated (child $\leq$4)}} \\
 & Post: not eligible & Post: eligible \\
$\quad$ Pre: not eligible & 84.3\% & 3.5\% \\
$\quad$ Pre: eligible & 2.6\% & 9.5\% \\
\midrule
\multicolumn{3}{l}{\textit{Control A (child 5--7)}} \\
 & Post: not eligible & Post: eligible \\
$\quad$ Pre: not eligible & 82.4\% & 3.7\% \\
$\quad$ Pre: eligible & 2.7\% & 11.1\% \\
\bottomrule\end{tabular}
\par\vspace{3pt}\footnotesize\raggedright \textit{Notes:} Each cell is the share of women in a given pre-to-post telework-eligibility transition, among matched women observed both before and after the second quarter of 2022 (15,048 treated and 6,587 control women). Off-diagonal cells are occupation switches. The treated and control distributions are similar, consistent with the null in Table~\ref{tab:mech_allocation}.
\end{table}


\begin{table}[htbp]
   \caption{\label{tab:triple_diff_outcomes} Triple Difference Across Outcomes}
   \centering
   \begin{tabular}{lcccccc}
      \tabularnewline \midrule \midrule
      Dependent Variables:                    & Home office   & Real income   & Usual hours & Employed  & In labor force & Maternity leave\\  
      Model:                                  & (1)           & (2)           & (3)         & (4)       & (5)            & (6)\\  
      \midrule
      \emph{Variables}\\
      Treated (child $\leq$4)                 & -0.004$^{**}$ & -40.2         & 0.040       & -0.0006   & 0.005$^{**}$   & 0.0002\\   
                                              & (0.002)       & (42.0)        & (0.118)     & (0.003)   & (0.002)        & (0.0002)\\   
      Treated $\times$ Female                 & 0.007$^{***}$ & 18.2          & -0.287      & -0.009    & -0.012$^{**}$  & 0.013$^{***}$\\   
                                              & (0.003)       & (44.5)        & (0.181)     & (0.006)   & (0.005)        & (0.002)\\   
      Treated $\times$ Post                   & 0.002         & 161.8$^{***}$ & 0.129       & 0.005     & -0.0006        & -0.0002\\   
                                              & (0.002)       & (58.0)        & (0.148)     & (0.004)   & (0.003)        & (0.0002)\\   
      Treated $\times$ Post $\times$ Female   & -0.003        & -129.6$^{**}$ & -0.054      & 0.005     & 0.011$^{**}$   & -0.008$^{***}$\\   
                                              & (0.004)       & (60.7)        & (0.237)     & (0.006)   & (0.006)        & (0.002)\\   
      \midrule
      \emph{Fixed-effects}\\
      Individual                              & Yes           & Yes           & Yes         & Yes       & Yes            & Yes\\  
      Female-Year-quarter                     & Yes           & Yes           & Yes         & Yes       & Yes            & Yes\\  
      \midrule
      \emph{Fit statistics}\\
      Observations                            & 1,819,196     & 1,202,578     & 1,228,429   & 1,819,196 & 1,819,196      & 1,819,196\\  
      R$^2$                                   & 0.666         & 0.887         & 0.759       & 0.802     & 0.795          & 0.546\\  
      \midrule \midrule
      \multicolumn{7}{l}{\emph{Clustered (Household) standard-errors in parentheses}}\\
      \multicolumn{7}{l}{\emph{Signif. Codes: ***: 0.01, **: 0.05, *: 0.1}}\\
   \end{tabular}
   
   \par \raggedright 
   Each column is a triple-difference regression with the same specification as the first-stage triple difference. Reported: the men effect (Treated $\times$ Post) and the female differential (Treated $\times$ Post $\times$ Female). $^{*}$/$^{**}$/$^{***}$: 10/5/1\%.
\end{table}




\begin{figure}[H]
\centering
\includegraphics[width=0.8\textwidth]{fig08_control_window_sweep.pdf}
\caption{Control-window robustness. Each point is the first-stage difference-in-differences estimate as the comparison group is widened from a youngest child aged 5--6 up to 5--12; bars are 95\% confidence intervals. The estimate is stable across windows.}
\label{fig:window}
\end{figure}
